\section{线性基}

用于维护一组数的线性基，可求子集最大异或和、判断某数能否由集合内异或表示、求基的秩等。\texttt{template<int N>} 为最高可处理的二进制位（默认 30，可支持 $\le 2^{31}-1$ 的整数）。

\texttt{insert(x)} 插入 \texttt{x}，返回 \texttt{true} 表示新增了独立基向量，\texttt{false} 表示 \texttt{x} 可由现有基表示。\texttt{operator[](i)} 取第 \texttt{i} 位的基向量，\texttt{count()} 返回基的秩，\texttt{operator==} 判断两个基是否等价（可表示的异或空间相同）。

\begin{lstlisting}
template<int N = 30>
struct Basis{
    using i64 = long long;
private:
    std::array<i64,N+1> b{};
    int cnt = 0;
public:
    Basis() {
        std::fill(b.begin(),b.end(),0);
    }
    bool insert(i64 x) {
        if (!x) return 0;
        for (int i = N;i >= 0;i--){
            if (!(x >> i & 1)) continue;
            if (!b[i]) {
                b[i] = x;
                for (int j = N;j >= 0;j--){
                    if (i == j || !b[j]) continue;
                    if (b[j] >> i & 1) b[j] ^= x;
                }
                ++cnt;
                return 1;
            } else x ^= b[i];
        }
        return 0;
    }
    i64 operator[](const int& i){return b[i];}
    i64 operator[](const int& i) const {return b[i];}
    bool operator==(const Basis& rhs) const {
        if (cnt != rhs.cnt) return 0;
        Basis t = *this;
        for (int i = N;i >= 0;i--)
            if (rhs.b[i]) t.insert(rhs.b[i]);
        return t.cnt == cnt;
    }
    int count() const {
        return cnt;
    }
};
\end{lstlisting}
